\documentclass{beamer}
\usetheme{CambridgeUS}
\usepackage[utf8]{inputenc}
\usepackage[russian]{babel}
\usepackage[T1]{fontenc}
\usepackage{amssymb,amsmath,graphicx,indentfirst}
\usepackage{enumerate}
\usepackage{fixltx2e}

\usepackage{caption}
\captionsetup[figure]{labelformat=empty}

\usepackage{listings}
\lstset{keywordstyle=\color{blue}\bfseries}
\lstset{basicstyle=\scriptsize}
\lstset{extendedchars=false, language=Caml, defaultdialect=[Objective]Caml}

\author{Иван Веселов}
\institute{Barclays Capital}
\date{28 октября 2011 г.}

\begin{document}

\begin{frame}{План лекции}
\begin{itemize}
\item
  Лирическое отступление: финансы и ФП
\item
  Свёртки на синтаксических деревьях
\item
  Обобщение свёрток: катаморфизмы
\item
  Почему ФП?
\end{itemize}

\end{frame}\begin{frame}{Barclays Capital}

\begin{figure}[htbp]
\centering
\includegraphics{lecture6/barclays.eps}
\caption{}
\end{figure}

\begin{itemize}
\item
  Barclays Capital -- инвестиционный банк
\item
  Предоставляет брокерские услуги, анализ рынков и т.д.
\item
  Наш проект: создание формального языка для описания деривативов
\end{itemize}
\end{frame}\begin{frame}{Пример дериватива}

\begin{itemize}
\item
  У Петра имеется 100 акций компании Apple
\item
  Цена 1 акции сегодня: \$5
\item
  Он может продать их сейчас или продать позже
\item
  Пётр любит Apple и надеется на то, что акции подорожают, но не слишком
  уверен в этом
\item
  Как ему застраховаться от возможных убытков?
\item
  Купить у нас дериватив!
\end{itemize}
\end{frame}\begin{frame}{Пример дериватива}

\begin{itemize}
\item
  Банк может заключить с Петром контракт, что

  \begin{enumerate}[1.]
  \item
    eсли акции подешевеют -- банк выплатит разницу в цене
  \item
    в противном случае -- банк ничего не выплачивает
  \end{enumerate}
\item
  Математически:

  \texttt{max(S0 - S1, 0)}
\item
  Пётр может воспользоваться этим контрактом через год.
\item
  Таким образом, если через год цена 1 акции составит 2 доллара, Пётр
  будет застрахован от убытков и сможет продать акции по той же цене,
  что и в день заключения контракта.
\item
  Однако за эту услугу ``страховки'' Петру придётся заплатить какую-то
  сумму, чтобы подписать такой контракт.
\item
  Дериватив -- это и есть контракт такого рода, зависящий от цены
  базового актива (цены акций, валют и т.д.).
\end{itemize}
\end{frame}\begin{frame}{Опционы}

\begin{itemize}
\item
  Базовые активы -- это что-то более конкретное (цена валюты раньше
  соответствовала ``золотому запасу'', акция отражает заинтересованность
  во владении часть компании)
\item
  Деривативы -- более абстрактное, это просто функция от базовых активов.
  Цена деривативов является производной от цен активов (отсюда
  название).
\item
  Дериватив, предложенный Петру, называется ``опцион''
\item
  Опционы могут быть сложнее (``экзотические опционы''), например:
\item
  Вычислить среднее значение цены за последний год и выплатить его
\item
  Сравнить цены нескольких различных акций и выплатить минимум
\item
  Выплатить взвешенное среднее цен нескольких акций
\end{itemize}
\end{frame}

\begin{frame}[fragile]
\frametitle{Пример в коде}

\begin{itemize}
\item
  Наш язык описывает такие контракты и выплаты по ним в виде, очень
  похожем на язык программирования.
\item
  Функция для описания выплаты Петру
\end{itemize}

\begin{lstlisting}[language=Haskell]
payout :: Market -> Double
payout market = 
  let s0 = observe market "28-10-2011" "AAPL"
      s1 = observe market "28-10-2012" "AAPL"
  in max (s0 - s1) 0
\end{lstlisting}
\end{frame}

\begin{frame}{Описание контрактов в коде}

\begin{itemize}
\item
  В чём преимущество описания в виде программы?

  \begin{enumerate}[1.]
  \item
    Гораздо лучше бумажного документа со словесным описанием, который
    может допускать неоднозначности
  \item
    Поскольку это программа, её можно выполнить и точно посчитать
    выплату
  \end{enumerate}
\item
  Недостатки

  \begin{itemize}
  \item
    Выполнить можно только когда наступит дата выплаты (цены акций -
    случайные величины)
  \end{itemize}
\end{itemize}
\end{frame}

\begin{frame}{Определение стоимости опциона}

\begin{itemize}
\item
  Однако, для того, чтобы продавать опционы нужно оценивать выплату
  заранее.
\item
  Для того, чтобы оценить выплату ещё до его продажи нужно примерно
  знать какими будут значения случайных величин.
\item
  Одно из решений: симуляция случайных величин методом Монте-Карло.

  \begin{itemize}
  \item
    каждая случайная величина имеет свою модель
  \item
    этим занимается отдел аналитиков
  \end{itemize}
\item
  Мы можем запустить 100000 раз симуляцию со случайной ценой акции,
  распределённой согласно своей модели, для каждого раза посчитать
  выплату, усреднить и попросить у Петра немного больше -- за сервис
\item
  Теоретически, никто не мешает Петру проделать это самому, чтобы
  оценить риск, но как правило у него нет квалификации или нужных
  инструментов
\end{itemize}
\end{frame}\begin{frame}{Инструменты}

\begin{itemize}
\item
  У нас такие инструменты есть!
\item
  FPF. Functional Payout Framework. Написан на Haskell.
\item
  Включает в себя язык описания деривативов.
\item
  Выполняет такие функции:
\item
  Генерация кода для оценки выплат (pricing instructions) по контрактам
\item
  Генерация описаний контрактов в виде формул
\item
  Различные анализы контрактов
\item
  Многое другое
\end{itemize}
\end{frame}

\begin{frame}{Пример описания контракта}

\begin{itemize}
\item
  Код скриптов представляет собой код на Haskell
\item
  Финансисты не очень хорошо читают Haskell, им понятнее форумы, потому
  мы генерируем описания контрактов в TeX
\item
  См. отдельный слайд
\end{itemize}
\end{frame}\begin{frame}{Извлечение информации}

\begin{itemize}
\item
  Важная информация, которую мы можем получить из контракта-скрипта:
\item
  Даты наблюдения
\item
  Множество базовых активов дериватива

  \begin{itemize}
  \item
    например названия задействованных в контракте акций
  \end{itemize}
\item
  Точки разрыва функции выплаты

  \begin{itemize}
  \item
    чтобы вовремя на них реагировать
  \end{itemize}
\end{itemize}
\end{frame}\begin{frame}{Как мы можем получить эту информацию?}

\begin{itemize}
\item
  Синтаксический разбор скрипта
\item
  Построение синтаксического дерева
\item
  Анализ дерева с помощью свёрток
\end{itemize}
\end{frame}

\begin{frame}{Cинтаксические деревья}
\begin{itemize}
\item
  \textbf{Синтаксическое дерево} -- это промежуточное представление кода
  программы в компиляторах языков программирования.
\item
  В процессе проверки синтаксического разбора программы парсер строит
  синтаксическое дерево, которое потом подвергается различным
  манипуляциям для оптимизации, проверки типов и т.д.
\item
  Вспомните лабораторную по синтаксическим деревьям регулярных
  выражений.
\end{itemize}
\end{frame}\begin{frame}{Свёртки для алгебраических типов данных}

\begin{itemize}
\item
  Суть \emph{свёртки} определяется названием: она \textbf{сворачивает}
  сложную структуру данных в некоторое значение. Так можно решить много
  разных задач, в том числе и поставленные выше.
\item
  Изобразим на доске пример сложной структуры в виде дерева, где каждый
  узел -- это конструктор.
\item
  Рядом изобразим дерево аналогичной структуры, только с операциями в
  узлах. Каждому конструктору соответствует отдельная операция.
\end{itemize}
\end{frame}

\begin{frame}[fragile]
\frametitle{Cинтаксические деревья}
\begin{itemize}
\item
  Мы будем рассматривать упрощённую версию дерева из языка,
  определённого в FPF.
\item
  Тип такого дерева:
\end{itemize}

\begin{lstlisting}
type expr =
    | Const of int
    | Var of string
    | Add of expr * expr
    | Mult of expr * expr
    | Observe of date * string
\end{lstlisting}

\end{frame}

\begin{frame}[fragile]
\frametitle{Примеры выражений}

\begin{itemize}
\item $observe "28-10-2011" "AAPL"$
\end{itemize}
\begin{lstlisting}
  Observe(Date "28-10-2011", "AAPL")
\end{lstlisting}
\begin{itemize}
\item $2 + 3$
\end{itemize}
\begin{lstlisting}
  Add(Const 2, Const 3)
\end{lstlisting}
\begin{itemize}
\item $s0 + s1$
\end{itemize}
\begin{lstlisting}
  Add(Var "s0", Var "s1")
\end{lstlisting}
\begin{itemize}
\item $0 + (1 * x))$
\end{itemize}
\begin{lstlisting}
  Add(Const 0, Mult(Const 1, Var "x"))
\end{lstlisting}
\end{frame}

\begin{frame}{Задача}
\begin{itemize}
\item
  Написать функцию для расчёта выплаты, используя дерево.

  \begin{itemize}
  \item
    на входе: синтаксическое дерево контракта (выражение)
  \item
    на выходе: выплата (число)
  \end{itemize}
\item
  Это -- задача свёртки, по сути нам нужно свернуть дерево по неким
  правилам в одно значение.
\item
  Теперь нужно определить операции, соответствующие разным конструкторам
  типа, которые будут решать задачу.
\end{itemize}
\end{frame}\begin{frame}{Свёртки и конструкторы}

\begin{itemize}
\item
  Вспомним о свёртках на списках, чтобы уяснить важный момент
\item
  Какие аргументы у функции свёртки для списков?

  \begin{enumerate}[1.]
  \item
    Значение на которое мы заменяем пустой список.
  \item
    Функция для комбинирования элемента списка с результатом свёртки
    хвоста списка.
  \end{enumerate}
\item
  Сравните вычисление правой свёртки на списке и структуру самого
  списка:
\end{itemize}

\begin{equation*}
  x_1 \oplus (x_2 \oplus (x_3 \oplus (\ldots \oplus (x_n \oplus base))))
\end{equation*}
\begin{equation*}
  x_1 :: (x_2 :: (x_3 :: (\ldots :: (x_n :: []))))
\end{equation*}

\begin{itemize}
\item
  Видно, что правая свёртка это по сути замена конструктора списка
  \texttt{(::)} на некую операцию. При этом финальный пустой список
  заменяется на некое значение \texttt{base}.
\end{itemize}
\end{frame}

\begin{frame}[fragile]
\frametitle{Свёртки и конструкторы}

\begin{itemize}
\item
  Список -- это алгебраический тип данных с двумя конструкторами
\end{itemize}
\begin{lstlisting}
type 'a list =
  | Nil
  | Cons of 'a * ('a list)
\end{lstlisting}
\begin{itemize}
\item
  Первый аргумент \texttt{fold} соответствует пустому списку
  (конструктор \texttt{Nil})
\item
  Второй аргумент соответствует голове и свёртке хвоста (конструктор
  \texttt{Cons})
\item
  Важно: свёртка -- это замена конструкторов типа на заданные функции
\item
  Таким образом теперь нам надо заменять конструкторы дерева на некие
  функции
\end{itemize}
\end{frame}

\begin{frame}{Преобразование типа данных в ``контекст''}
\begin{itemize}
\item
  Можно решить данную задачу с помощью рекурсивно определённой функции.
  Но можно сделать это более элегантно, задав простые и нерекурсивные
  способы комбинирования значений, а всю рекурсию спрятать в оператор
  свёртки.
\item
  Введём понятия ``контекста свёртки''. В контексте будут храниться уже
  вычисленные значения предыдущих свёрток.
\item
  Чтобы получить тип контекста применим преобразование к типу
  \texttt{expr}, чтобы получить новый тип данных.
\end{itemize}
\end{frame}

\begin{frame}[fragile]
\frametitle{Преобразование}

Было:
\begin{lstlisting}
type expr =
    | Const of int
    | Var of string
    | Add of expr * expr
    | Mult of expr * expr
    | Observe of date * string
\end{lstlisting}
\begin{enumerate}[1.]
\item
  Вводим новый параметр типа \texttt{a}
\item
  Зaменяем рекурсивные вхождения expr в оригинальном типе на параметр
  \texttt{a}
\item
  Переименовываем конструкторы, добавляя \texttt{Context}.
\end{enumerate}
Стало:
\begin{lstlisting}
type 'a exprContext =
    | ConstContext of int
    | VarContext of string
    | AddContext of 'a * 'a
    | MultContext of 'a * 'a
    | ObserveContext of date * string
\end{lstlisting}
\end{frame}

\begin{frame}{Функции на контексте}

\begin{itemize}
\item
  Наш контекст -- типизированный, в нём могут храниться определённые
  значения типа \texttt{а}, а именно -- промежуточные результаты.
\item
  Теперь определим функцию на контексте, определяющую свёртку
\item
  Она будет иметь тип \texttt{'a exprContext -\textgreater{} a}. Эта
  функции будет нерекурсивной и её достаточно для того, чтобы полностью
  определить свёртку.
\item
  Для нашей задачи результатом свёртки будет \texttt{int}, т.к.
  результат вычисления выражения -- число, а именно результат определяет
  тип \texttt{а} в контексте.
\end{itemize}
\end{frame}

\begin{frame}[fragile]
\frametitle{Функции на контексте}

\begin{itemize}
\item
  Таким образом, нам надо задать функцию вида
  \texttt{'int exprContext -\textgreater{} int}
\end{itemize}
\begin{lstlisting}
// eval : 'int exprContext -> int 
let eval expr = 
    match expr with
    | ConstContext x        -> x
    | VarContext s          -> 42 // lookup in env 
    | AddContext (x, y)     -> x + y
    | MultContext (x, y)    -> x * y
    | ObserveContext (d, a) -> lookup d a market
\end{lstlisting}
\begin{itemize}
\item
  Здесь мы не делаем рекурсивных вызовов, а просто описываем как
  комбинировать элементы, считая что предыдущие вызовы у нас уже
  вычислены. Можно сказать, что эта функция знает как свернуть один
  уровень структуры данных.
\item
  Изобразим на доске картинку, которая сравнит функцию на деревьях и
  функцию на контекстах.
\item
  Функция \texttt{eval} нерекурсивна!
\end{itemize}
\end{frame}

\begin{frame}[fragile]
\frametitle{Катаморфизм}

\begin{itemize}
\item
  Как теперь ``запустить'' эту функцию, чтобы она приняла дерево и
  выдала результат?
\item
  Для этого нам нужно преобразовать функцию \texttt{f}, которая
  определяет свёртку, в реальную свёртку. То есть у нас есть определение
  того, как свернуть один уровень дерева, и нам нужно свернуть всё
  дерево, обладая только этой функцией и знанием о типе дерева. Это
  делается с помощью функции такого рода:
\end{itemize}
\begin{lstlisting}
// toFold : ('a exprContext -> 'a) -> expr -> 'a
let rec toFold ctx f = 
    match ctx with
    | Const x        -> f (ConstContext x)
    | Var   s        -> f (VarContext s)
    | Add  (x, y)    -> f (AddContext (toFold f x, toFold f y))
    | Mult (x, y)    -> f (MultContext (toFold f x, toFold f y))
    | Observe (d, a) -> f (ObserveContext(d, a))
\end{lstlisting}
\end{frame}

\begin{frame}{Катаморфизм}

\begin{itemize}
\item
  Именно в этой функции ``зашита'' рекурсия, однако прелесть в том, что
  эта функция может быть обобщена на любой рекурсивный тип данных
  (практически любой) и её не нужно будет создавать заново.
\item
  Такое обобщение свёрток на алгебраические структуры данных называется
  ``катаморфизмом''.
\item
  ``ката'' -- по-гречески означает ``вниз''. Однокоренное слово
  ``катастрофа'', то есть ``разрушение''. Катаморфизм как бы дробит,
  разрушает структуру данных, преобразовывая её в какое-то единое
  значение.
\item
  Можно переименовать функцию \texttt{toFold} в \texttt{cata}
\end{itemize}
\end{frame}

\begin{frame}[fragile]
\frametitle{Решение задачи про вычислитель}

\begin{itemize}
\item
  Соответственно, финальная функция, которая будет решать нашу задачу
  выглядит так:
\end{itemize}
\begin{lstlisting}
// e = 2 + 3 * 4
let e = Add(Const 2, Mult(Const 3, Const 4))
let solve e = cata eval e
let main = printfn "%A" (solve e)
// output: 14
\end{lstlisting}
\begin{itemize}
\item
  Важно, что \texttt{f} -- нерекурсивна, вся рекурсия спрятана в
  \texttt{toFold}.
\end{itemize}
\end{frame}

\begin{frame}{Задача оптимизации выражений}

\begin{itemize}
\item
  Предположим нам нужно написать оптимизатор выражений.
\item
  К примеру у нас есть выражение \texttt{1 * (x + 0)}. Вычислить мы его
  не можем, поскольку неизвестен \texttt{x} однако мы можем
  оптимизировать, выбросив умножение на 1 и сложение с нулём, т.к. они
  не меняют результат.
\end{itemize}
\end{frame}

\begin{frame}[fragile]
\frametitle{Решение}

\begin{itemize}
\item
  Для её решения определяем ещё одну функцию для контекста.
\item
  Теперь параметром \texttt{a} в контексте будет не число, а
  \texttt{expr}, поскольку результатом нашей функции является выражение.
\item
  Функция будет задавать наши правила оптимизации.
\end{itemize}
\begin{lstlisting}
// optimizer : expr exprContext -> expr
let optimizer e =
  match e with  
    | ConstContext x           -> Const x
    | VarContext s             -> Var s
    | AddContext (Const 0, y)  -> y
    | AddContext (y, Const 0)  -> y
    | AddContext (x, y)        -> Add (x, y)
    | MultContext (Const 1, y) -> y
    | MultContext (y, Const 1) -> y
    | MultContext (x, y)       -> Mult (x, y)
    | ObserveContext (d, a)    -> Observe (d, a)
// expr for '1 * (x + 0)'
let e2 = Mult(Const 1, Add(Var "x", Const 0))
let main = printfn "%A" (toFold optimizer e2)
// output: Var "x"
\end{lstlisting}
\end{frame}

\begin{frame}[fragile]
\frametitle{Катаморфизм на списках}

\begin{itemize}
\item
  Можно провести аналогичную операцию со введением контекста для
  списков, чтобы показать как будет выглядеть свёртка в варианте с
  контекстом.

\begin{lstlisting}
// type for list fold context
type 'a listContext =
  | NilContext
  | ConsContext of int * 'a

// function which matters (defines combination,
// one case for one constructor)
sumC NilContext        = 0
sumC (ConsContext x s) = x + s

// sum as catamorphism
sum = cata sumC
\end{lstlisting}
\end{itemize}
\end{frame}

\begin{frame}[fragile]
\frametitle{Запись свёртки на деревьях через foldExpr}

\begin{itemize}
\item
  Ещё один вариант определения свёрток на синтаксических деревьях:
  написание функции, аналогичной функции fold для списков.
\end{itemize}
\begin{lstlisting}
// foldExpr : (a->a->a) -> (a->a->a) -> (int->a) -> 
// -> (string->a) -> (date->string->a) -> expr -> a
let foldExpr add mult constf var obs e =
    let rec go e = 
      match e with 
        | Add (e1, e2)   -> add (go e1) (go e2)
        | Mult (e1, e2)  -> mult (go e1) (go e2)
        | Const c        -> constf c
        | Var name       -> var name
        | Observe (d, a) -> obs d a
    go e
\end{lstlisting}
\begin{itemize}
\item
  Здесь, как и в случае с \texttt{fold} на списках, мы задаём по одной
  функции для каждого из конструкторов.
\item
  Тип функции foldExpr довольно ужасен (много сложных параметров)
\end{itemize}
\end{frame}

\begin{frame}[fragile]
\frametitle{Ещё один пример вычислителя}

Тогда вычислитель будет выглядеть так:

\begin{lstlisting}
// yet another evaluator
let eval2 env market e =
    let add x y  = x + y
    let mult x y = x * y
    let constF c = c
    let var s    = lookup s env
    let obs d a  = lookupMarket d a market
    foldExpr add mult constF var obs e

// e = 2 + 3 * 4
let e = Add(Const 2, Mult(Const 3, Const 4))
let main = printfn "%A" (eval2 emptyEnv getMarketForToday e)
// output: 14
\end{lstlisting}
\end{frame}

\begin{frame}{Сравнение \texttt{foldl} на списках и \texttt{cata}}

\begin{itemize}
\item
  Зачем это нужно? В чём преимущество использования контекста перед
  обычным использованием \texttt{foldr/foldl} для списков?

  \begin{enumerate}[1.]
  \item
    В функции \texttt{foldr/l} каждому типу конструктора соответствовал
    свой аргумент (базовый элемент для пустого списка, комбинирующая
    функция -- для непустого), таким образом количество аргументов
    соответствует количеству конструкторов. Во многих случаях это
    неудобно, т.к. конструкторов может быть очень много (например в
    реальном синтаксическом дереве). Использовать функции с десятками
    аргументов непрактично: это нечитабельно, сложно поддерживать, нужно
    менять функцию каждый раз при добавлении нового конструктора.
  \item
    При создании нового типа данных для него пришлось бы создавать свою
    функцию \texttt{foldX}, которая задавала бы определённую схему
    рекурсивной свёртки. В нашем же случае всё, что нам нужно -- это
    задать тип данных для контекста, всё остальное может быть выведено
    автоматически с помощью ``комбинатора неподвижной точки'' над типами
    данных. Как именно -- будет рассмотрено в следующем семестре. Таким
    образом, подход намного более универсальный.
  \end{enumerate}
\item
  Общее достоинство обоих методов в том, что конечному пользователю
  функций не нужно явно задавать рекурсивные функции. Гораздо проще
  определять нерекурсивные функции, а рекурсию выполнять с помощью
  оператора (\texttt{fold} или \texttt{cata}). Таким образом схема
  рекурсии чётко определена и не будет содержать неожиданностей.
\end{itemize}
\end{frame}\begin{frame}{Barclays Capital и ФП}

Почему функциональное программирование и Haskell?

\begin{itemize}
\item
  Хорошо соответствует предметной области: выплаты по деривативам
  представляются в виде математических функций (выходы чётко
  соответствуют входам, нет сайд-эффектов): их можно сочетать,
  передавать в функции высшего порядка.
\item
  Легко сделать поддержку функций высшего порядка в скриптах. Это
  удобнее для скриптеров, позволяет более кратко описывать важные
  концепции.
\item
  Сопоставление с образцом (pattern matching) удобно для преобразования
  деревьев. Операции на деревьях -- одни из основных в компиляторах и
  всяких преобразователях или генераторах кода.
\item
  Статическая типизация и отсутствие сайд-эффектов позволяет писать
  ``код, в который можно верить''.
\item
  Более гибкий код, иногда даже слишком гибкий! Слабое связывание
  компонентов -- легко делать рефакторинги.
\end{itemize}
\end{frame}\begin{frame}{Приходите к нам!}

\begin{itemize}
\item
  Возможен internship (летняя практика) для талантливых студентов.
\item
  Есть вакансии для Haskell-программистов и скриптеров.
\item
  Знание Haskell желательно, но не обязательно, опыт с F\# подходит.
\end{itemize}
\end{frame}\begin{frame}{Задачи}

Несколько простых задач.

\begin{enumerate}[1.]
\item
  Реализовать катаморфизм на синтаксическом дереве, который возвращает
  множество переменных, использованных в выражении.
\item
  Разработать алгебраический тип и свёртки для следующих задач.

  2.1. У нас есть дерево документа, к примеру этой презентации.
  Элементами дерева могут быть:

  \begin{itemize}
  \item
    параграф
  \item
    картинка
  \item
    ссылка
  \item
    текст
  \end{itemize}
  В каждом параграфе может быть произвольное количество картинок,
  подпараграфов, текста, ссылок и т.д.

  Задача состоит в том, чтобы извлечь все ссылки для того, чтобы вывести
  их в конце презентации. Напишите для этого катаморфизм.

  2.2. Напишите катаморфизм, который удалит все картинки из документа.
\end{enumerate}
\end{frame}

\begin{frame}{Контакты}

\begin{itemize}
\item Иван Веселов
\item Ivan.Veselov@BarclaysCapital.com
\item veselov@gmail.com
\end{itemize}
\end{frame}\begin{frame}{Полезные ссылки}

\begin{enumerate}[1.]
\item
  ``Commercial Uses: Going functional on exotic trade'' -- статья про FPF

  http://arbitrary.name/papers/fpf.pdf
\item
  Набор статей о катаморфизмах в F\# от одного из ведущих программистов
  команды F\#

  http://lorgonblog.wordpress.com/2008/04/05/catamorphisms-part-one/
\item
  ``Functional Programming with Bananas, Lenses, Envelopes and Barbed
  Wire''

  Одна из первых статей про морфизмы в контексте программирования
  (сложная)

  http://citeseer.ist.psu.edu/viewdoc/summary?doi=10.1.1.41.125
\end{enumerate}

\end{frame}\begin{frame}{Спасибо за внимание!}
  \begin{center}
    Задавайте вопросы.
  \end{center}
\end{frame}

\end{document}
